% !TEX program = xelatex
\documentclass[a4paper,10pt]{article}
\usepackage[UTF8,fontset=none]{ctex}
\usepackage{xeCJK}
\setCJKmainfont{Songti SC}
\setCJKsansfont{STHeiti}
\setCJKmonofont{Songti SC}
\usepackage[margin=2.2cm]{geometry}
\usepackage{tikz}
\usetikzlibrary{shapes.geometric, arrows.meta, positioning, calc, fit}
\usepackage{amsmath, amssymb}
\usepackage{longtable}
\usepackage{array}
\usepackage{booktabs}
\usepackage{multirow}
\usepackage{enumitem}
\usepackage{fancyhdr}
\usepackage{titlesec}
\usepackage{hyperref}
\usepackage{xcolor}
\usepackage{tabularx}

\hypersetup{
  colorlinks=true,
  linkcolor=blue!60!black,
  urlcolor=blue!60!black,
  bookmarksnumbered=true
}

% 页眉页脚
\pagestyle{fancy}
\fancyhf{}
\fancyhead[L]{\small clampValue~---~函数设计文档}
\fancyhead[R]{\small 内部公开}
\fancyfoot[C]{\thepage}
\renewcommand{\headrulewidth}{0.4pt}

% 层级编号
\setcounter{secnumdepth}{3}
\setcounter{tocdepth}{3}

% 表格列宽辅助
\newcolumntype{L}[1]{>{\raggedright\arraybackslash}p{#1}}
\newcolumntype{C}[1]{>{\centering\arraybackslash}p{#1}}

% ============================================================
\begin{document}

% ==================== 封面表格 ====================
\thispagestyle{empty}

\begin{center}
{\CJKfamily{zhsong}\bfseries\Large 函数设计文档}
\end{center}

\vspace{8pt}

\begin{center}
\renewcommand{\arraystretch}{1.5}
\begin{tabular}{|L{3.2cm}|L{4.5cm}|L{3.2cm}|L{4.5cm}|}
\hline
\textbf{函数英文名} & clampValue &
\textbf{函数中文名} & 数值限幅截断 \\
\hline
\textbf{函数类型} & 元件 (Element) &
\textbf{版本} & 2.0 \\
\hline
\textbf{作者} & 许庆 &
\textbf{日期} & 2026-06-26 \\
\hline
\textbf{联系方式} & \multicolumn{3}{l|}{18612426814~/~qingxu@tsinghua.edu.cn} \\
\hline
\textbf{密级} & 内部公开 &
\textbf{AI辅助} & Cursor Agent (Claude Opus 4.6) \\
\hline
\end{tabular}
\end{center}

\vspace{12pt}

\begin{center}
\renewcommand{\arraystretch}{1.4}
\begin{tabular}{|C{1.2cm}|C{2cm}|L{5cm}|C{2cm}|C{2cm}|}
\hline
\multicolumn{5}{|c|}{\textbf{更改记录}} \\
\hline
\textbf{版本} & \textbf{日期} & \textbf{更改说明} & \textbf{作者} & \textbf{辅助AI} \\
\hline
2.0 & 2026-06-26 & 初始版本 & 许庆 & Claude Opus 4.6 \\
\hline
\end{tabular}
\end{center}

\newpage

% ==================== 目录 ====================
\tableofcontents
\newpage

% ============================================================
%  1. 函数需求
% ============================================================
\section{函数需求}

\subsection{功能概述}

\texttt{clampValue} 将输入数值截断到指定闭区间 $[\text{minVal}, \text{maxVal}]$ 内。
若输入值低于下限则返回下限，高于上限则返回上限，否则返回原值。

\subsection{输入/输出/返回值}

\renewcommand{\arraystretch}{1.3}
\begin{longtable}{|C{1.5cm}|L{3.5cm}|C{2cm}|L{8cm}|}
\hline
\textbf{类别} & \textbf{名称} & \textbf{类型} & \textbf{说明} \\
\hline
\endfirsthead
\hline
\textbf{类别} & \textbf{名称} & \textbf{类型} & \textbf{说明} \\
\hline
\endhead
输入 & \texttt{value} & double & 待截断数值 \\
\hline
输入 & \texttt{minVal} & double & 截断下限 \\
\hline
输入 & \texttt{maxVal} & double & 截断上限 \\
\hline
返回值 & --- & double & 截断后的数值，$\in [\text{minVal}, \text{maxVal}]$ \\
\hline
\end{longtable}

% ============================================================
%  1.2 算法设计
% ============================================================
\subsection{算法设计}

数值限幅截断的数学表达式为：
\begin{equation}
  \text{output} = \max\!\big(\text{minVal},\; \min(\text{maxVal},\; \text{value})\big)
\end{equation}

等价的分段函数形式：
\begin{equation}
  \text{output} =
  \begin{cases}
    \text{minVal}, & \text{value} < \text{minVal} \\
    \text{maxVal}, & \text{value} > \text{maxVal} \\
    \text{value},  & \text{otherwise}
  \end{cases}
\end{equation}

实现中采用两步条件赋值避免调用 \texttt{std::min}/\texttt{std::max}：
\begin{enumerate}
  \item $\text{clampedLow} = (\text{value} < \text{minVal})\ ?\ \text{minVal} : \text{value}$
  \item $\text{result} = (\text{clampedLow} > \text{maxVal})\ ?\ \text{maxVal} : \text{clampedLow}$
\end{enumerate}

% ============================================================
%  1.3 程序流程图
% ============================================================
\subsection{程序流程图}

\begin{center}
\begin{tikzpicture}[
  node distance=1.0cm,
  startstop/.style={rectangle, rounded corners, minimum width=3.5cm, minimum height=0.7cm,
    text centered, draw=black, fill=blue!8, font=\small},
  process/.style={rectangle, minimum width=4.5cm, minimum height=0.7cm,
    text centered, draw=black, fill=orange!8, font=\small},
  decision/.style={diamond, aspect=2.5, minimum width=2cm,
    text centered, draw=black, fill=green!8, font=\small},
  arrow/.style={-{Stealth[length=2.5mm]}, thick}
]

\node (start) [startstop] {开始};
\node (input) [process, below=of start] {读入 value, minVal, maxVal};
\node (chklow) [decision, below=of input] {value $<$ minVal?};
\node (setlow) [process, right=3cm of chklow] {clampedLow = minVal};
\node (keepval) [process, below=of chklow, yshift=-0.5cm] {clampedLow = value};
\node (chkhigh) [decision, below=of keepval] {clampedLow $>$ maxVal?};
\node (sethigh) [process, right=3cm of chkhigh] {result = maxVal};
\node (keepcl) [process, below=of chkhigh, yshift=-0.5cm] {result = clampedLow};
\node (output) [process, below=of keepcl] {返回 result};
\node (stop) [startstop, below=of output] {结束};

\draw [arrow] (start) -- (input);
\draw [arrow] (input) -- (chklow);
\draw [arrow] (chklow) -- node[above]{\small 是} (setlow);
\draw [arrow] (chklow) -- node[right]{\small 否} (keepval);
\draw [arrow] (setlow) |- (chkhigh);
\draw [arrow] (keepval) -- (chkhigh);
\draw [arrow] (chkhigh) -- node[above]{\small 是} (sethigh);
\draw [arrow] (chkhigh) -- node[right]{\small 否} (keepcl);
\draw [arrow] (sethigh) |- (output);
\draw [arrow] (keepcl) -- (output);
\draw [arrow] (output) -- (stop);

\end{tikzpicture}
\end{center}

% ============================================================
%  1.4 函数接口
% ============================================================
\subsection{函数接口}

\begin{verbatim}
double clampValue(const double value, const double minVal, const double maxVal);
\end{verbatim}

% ============================================================
%  1.5 输入输出模块图
% ============================================================
\subsection{输入输出模块图}

\begin{center}
\begin{tikzpicture}[
  box/.style={rectangle, draw=black, thick, minimum width=3.0cm, minimum height=0.6cm,
    text centered, fill=yellow!10, font=\small},
  core/.style={rectangle, draw=black, very thick, minimum width=4cm, minimum height=1.8cm,
    text centered, fill=orange!10, font=\normalsize\bfseries, rounded corners=3pt},
  arrow/.style={-{Stealth[length=2.5mm]}, thick}
]

\node (core) [core] {clampValue};

\node (in1) [box, left=3cm of core, yshift=0.8cm] {value (double)};
\node (in2) [box, left=3cm of core, yshift=0cm] {minVal (double)};
\node (in3) [box, left=3cm of core, yshift=-0.8cm] {maxVal (double)};

\node (out1) [box, right=3cm of core] {result (double)};

\draw [arrow] (in1) -- (core);
\draw [arrow] (in2) -- (core);
\draw [arrow] (in3) -- (core);
\draw [arrow] (core) -- (out1);

\node [above=0.1cm of in1, font=\small\bfseries] {输入 (Input)};
\node [above=0.1cm of out1, font=\small\bfseries] {输出 (Output)};

\end{tikzpicture}
\end{center}

% ============================================================
%  1.6 函数诸元表
% ============================================================
\subsection{函数诸元表}

\renewcommand{\arraystretch}{1.3}
\begin{longtable}{|L{3.8cm}|L{11.5cm}|}
\hline
\textbf{属性} & \textbf{说明} \\
\hline
\endfirsthead
\hline
\textbf{属性} & \textbf{说明} \\
\hline
\endhead
函数英文名 & \texttt{clampValue} \\
\hline
函数中文名 & 数值限幅截断 \\
\hline
函数类型 & 元件 (Element) \\
\hline
版本 & 2.0 \\
\hline
源文件 & \texttt{src/cpp/clampValue/clampValue.cpp} \\
\hline
头文件 & \texttt{include/cpp/clampValue/clampValue.hpp} \\
\hline
命名空间 & \texttt{waypoint\_follow} \\
\hline
输入数量 & 3（value, minVal, maxVal） \\
\hline
输出数量 & 1（返回值 double） \\
\hline
参数数量 & 0 \\
\hline
状态数量 & 0（无状态，纯函数） \\
\hline
下级函数数量 & 0 \\
\hline
动态内存分配 & 无 \\
\hline
外部依赖 & 无 \\
\hline
算法类别 & 数值截断 / 安全限幅 \\
\hline
\end{longtable}

% ============================================================
%  1.7 函数架构
% ============================================================
\subsection{函数架构}

无下级函数调用。\texttt{clampValue} 为叶节点元件，仅包含条件判断与赋值操作。

% ============================================================
%  1.8 数据结构定义
% ============================================================
\subsection{数据结构定义}

无自定义数据结构。所有输入输出均为 \texttt{double} 标量。

% ============================================================
%  1.9 测试用例
% ============================================================
\subsection{测试用例}

\subsubsection{正常测试用例}

{\small
\begin{longtable}{|C{0.6cm}|L{2.2cm}|L{4.5cm}|L{3.5cm}|L{3.8cm}|}
\hline
\textbf{编号} & \textbf{场景} & \textbf{输入} & \textbf{预期输出} & \textbf{验证准则} \\
\hline
\endfirsthead
\hline
\textbf{编号} & \textbf{场景} & \textbf{输入} & \textbf{预期输出} & \textbf{验证准则} \\
\hline
\endhead
N-1 &
值在区间内 &
$v{=}5.0$, $l{=}0.0$, $h{=}10.0$ &
$r = 5.0$ &
$r = v$（原值返回） \\
\hline
N-2 &
值等于下限 &
$v{=}0.0$, $l{=}0.0$, $h{=}10.0$ &
$r = 0.0$ &
$r = l$（边界值） \\
\hline
N-3 &
值等于上限 &
$v{=}10.0$, $l{=}0.0$, $h{=}10.0$ &
$r = 10.0$ &
$r = h$（边界值） \\
\hline
N-4 &
值低于下限 &
$v{=}{-}3.0$, $l{=}0.0$, $h{=}10.0$ &
$r = 0.0$ &
$r = l$（下截断） \\
\hline
N-5 &
值高于上限 &
$v{=}15.0$, $l{=}0.0$, $h{=}10.0$ &
$r = 10.0$ &
$r = h$（上截断） \\
\hline
\end{longtable}
}

\vspace{6pt}
{\small
\begin{longtable}{|L{2.5cm}|L{13cm}|}
\hline
\multicolumn{2}{|c|}{\textbf{正常测试用例符号说明}} \\
\hline
\textbf{符号} & \textbf{含义} \\
\hline
$v$ & 输入值 value \\
\hline
$l$ & 截断下限 minVal \\
\hline
$h$ & 截断上限 maxVal \\
\hline
$r$ & 返回值 result \\
\hline
\end{longtable}
}

\subsubsection{异常测试用例}

{\small
\begin{longtable}{|C{0.6cm}|L{2.0cm}|L{4.2cm}|L{4.0cm}|L{3.8cm}|}
\hline
\textbf{编号} & \textbf{场景} & \textbf{输入} & \textbf{预期输出/行为} & \textbf{验证准则} \\
\hline
\endfirsthead
\hline
\textbf{编号} & \textbf{场景} & \textbf{输入} & \textbf{预期输出/行为} & \textbf{验证准则} \\
\hline
\endhead
A-1 &
value 为 NaN &
$v{=}\text{NaN}$, $l{=}0$, $h{=}10$ &
条件比较结果为 false，返回 NaN &
函数不崩溃，正常返回 \\
\hline
A-2 &
value 为 $+\infty$ &
$v{=}+\infty$, $l{=}0$, $h{=}10$ &
$r = 10.0$（上截断） &
$+\infty > h$ 成立，被截断 \\
\hline
A-3 &
value 为 $-\infty$ &
$v{=}-\infty$, $l{=}0$, $h{=}10$ &
$r = 0.0$（下截断） &
$-\infty < l$ 成立，被截断 \\
\hline
A-4 &
minVal 为 NaN &
$v{=}5$, $l{=}\text{NaN}$, $h{=}10$ &
NaN 比较为 false，可能返回 $v$ 或 NaN &
函数不崩溃 \\
\hline
A-5 &
maxVal 为 NaN &
$v{=}5$, $l{=}0$, $h{=}\text{NaN}$ &
NaN 比较为 false，可能返回 $v$ 或 NaN &
函数不崩溃 \\
\hline
A-6 &
minVal $>$ maxVal &
$v{=}5$, $l{=}10$, $h{=}0$ &
$r = 0$（先截下限得 10，再截上限得 0） &
退化区间行为确定 \\
\hline
A-7 &
全部为 $+\infty$ &
$v{=}+\infty$, $l{=}+\infty$, $h{=}+\infty$ &
$r = +\infty$ &
函数不崩溃 \\
\hline
A-8 &
minVal $=$ maxVal &
$v{=}5$, $l{=}3$, $h{=}3$ &
$r = 3.0$ &
零宽区间直接返回该值 \\
\hline
A-9 &
极大值输入 &
$v{=}1.7\!\times\!10^{308}$, $l{=}0$, $h{=}100$ &
$r = 100.0$（上截断） &
double 最大值被正确截断 \\
\hline
A-10 &
极小负值输入 &
$v{=}{-}1.7\!\times\!10^{308}$, $l{=}{-}100$, $h{=}0$ &
$r = -100.0$（下截断） &
double 最小值被正确截断 \\
\hline
\end{longtable}
}

\vspace{6pt}
{\small
\begin{longtable}{|L{2.5cm}|L{13cm}|}
\hline
\multicolumn{2}{|c|}{\textbf{异常测试用例符号说明}} \\
\hline
\textbf{符号} & \textbf{含义} \\
\hline
$v$ & 输入值 value \\
\hline
$l$ & 截断下限 minVal \\
\hline
$h$ & 截断上限 maxVal \\
\hline
$r$ & 返回值 result \\
\hline
NaN & IEEE 754 非数值 (Not a Number)，任何与 NaN 的比较返回 false \\
\hline
$+\infty$ & IEEE 754 正无穷大 \\
\hline
$-\infty$ & IEEE 754 负无穷大 \\
\hline
$1.7\!\times\!10^{308}$ & 接近 double 最大有限值 (DBL\_MAX) \\
\hline
\end{longtable}
}

\end{document}
